The folder \texttt{SeymourLean/} is a Lean~4 project using Mathlib (tag
\texttt{v4.35.0}) that formalizes the \emph{entire} proof of
\cref{thm:main}. The main formal statement is
\begin{center}\small
\texttt{theorem Seymour.seymour\_second\_neighbourhood (V : Type) [Fintype V]
[DecidableEq V] [Nonempty V] (adj : V $\to$ V $\to$ Prop) [DecidableRel adj]
(hor : IsOriented adj) : $\exists$ v, (N1 adj v).card $\le$ (N2 adj v).card}
\end{center}
where \texttt{IsOriented} means irreflexive and asymmetric, \texttt{N1 v} is
the set of out-neighbors, and \texttt{N2 v} is the set of $w\ne v$ with
$v\not\to w$ and $v\to u\to w$ for some $u$. \texttt{lake build} checks it with
no \texttt{sorry}, and \texttt{\#print axioms} reports only Lean's standard
axioms \texttt{propext}, \texttt{Classical.choice} and \texttt{Quot.sound}.
The file \texttt{SeymourLean/build.log} records a full build together with
the output of \texttt{lake env lean CheckAxioms.lean}, so this can be checked
without rebuilding. The files follow the paper:
\begin{itemize}
\item \texttt{Reduction.lean}: minimal counterexample, \cref{lem:strong} and
\cref{lem:subset-deficit};
\item \texttt{BlowUp.lean}: \cref{lem:blowup};
\item \texttt{Contradiction.lean}: \cref{prop:no-concave} from the pruning lemma;
\item \texttt{Konig.lean}, \texttt{Gamma.lean}: K\H{o}nig's theorem (derived from
Mathlib's Hall theorem) and \cref{cl:cover-to-pruning};
\item \texttt{LinAlg.lean}, \texttt{Generic.lean}, \texttt{MatchingCore.lean},
\texttt{MatchingBound.lean}: Sylvester's inequality, the kernel lemma
(\cref{lem:kernel}), $NL=AB$, \cref{cl:conflict-kernel},
\cref{prop:matching-bound}, and \cref{lem:pruning};
\item \texttt{Final.lean}: the unconditional theorem;
\item \texttt{Insights.lean}: the remarks of the insights pass, see below.
\end{itemize}
The remarks added in \cref{sec:story,sec:ledger,sec:pruning,sec:contradiction}
are formalized in \texttt{Insights.lean}: \texttt{F1\_F2\_comm} ($\Ff\Fs=\Fs\Ff$)
and \texttt{kronecker\_comm} (\cref{rem:kronecker});
\texttt{diag\_compatible\_iff\_oriented} and
\texttt{exists\_admissible\_iff\_oriented} (\cref{lem:admissible-iff-oriented});
\texttt{conflict\_swap}, \texttt{Zset\_swap}, \texttt{Hset\_swap}
(\cref{rem:reversal}); \texttt{gain\_identity} (the exact form of
\cref{cl:gain}) and \texttt{moved\_identity} (\cref{rem:move-repair});
and \texttt{step\_ledger}, \texttt{step\_value}: for an arbitrary relation
and any admissible pair, pruning the complements gives an admissible pair of
value at least $\val(P,Q)-\sum\kappa$ (\cref{rem:algorithm}), with no
concavity assumed. Not formalized: the Dulmage--Mendelsohn reading of the
kernel lemma (\cref{rem:DM}), the computer observations of
\cref{rem:combinatorial}, and the running-time statement of
\cref{rem:algorithm}. \texttt{CheckAxioms.lean} prints the axioms of
these theorems as well.
A few steps are argued differently in Lean. Generic coefficients are taken
as the variables themselves in the fraction field of a polynomial ring
rather than as a real point. Only $\rank L\ge\ell$ is used, obtained from
Hall's theorem and an exchange argument, instead of the full statement
$\rank L=\nu(\Gamma_L)$. The rank count goes through the three-term complex of \cref{sec:rank}
(\texttt{ComplexCount.lean}: \texttt{complex\_count} is the counting of
\cref{cl:kerA,cl:rank-inequality}, and \texttt{dimension\_count} is
\cref{cl:rank-inequality} for the matrices of the paper); the earlier
version of the project used Sylvester's inequality twice instead. Strong connectivity is replaced by ``every nonempty
proper set has an arc leaving it''. The formalization found no mathematical
error in the paper. It also shows that the pruning lemma holds for an
arbitrary relation and that orientation is used exactly at the places listed
in \cref{rem:hypotheses}.
